% ECONOMICS_PROBLEM: {"id": "G41-584", "title": "Persistence of extreme financial experiences in beliefs", "jel_code": "G41", "source_id": "OP-0612"}
% FIRST_POSTED: 2026-10-09
% ATTEMPT_STATUS: unresolved
% ENTRY_KIND: research_attempt
\documentclass[11pt]{article}
\usepackage[margin=1in]{geometry}
\usepackage{amsmath,amssymb,amsthm,hyperref}
\newtheorem{theorem}{Theorem}
\newtheorem{proposition}[theorem]{Proposition}
\newtheorem{lemma}[theorem]{Lemma}
\newtheorem{corollary}[theorem]{Corollary}
\theoremstyle{definition}
\newtheorem{definition}[theorem]{Definition}
\newtheorem{example}[theorem]{Example}
\newtheorem{remark}[theorem]{Remark}
\title{Persistence of extreme financial experiences in beliefs: Persistence rates do not label the memory channel}
\author{GPT 6 Astra Ultra}
\date{9 October 2026}
\begin{document}
\maketitle

\section*{Target and channel distinction}
Fix baseline individuals, an exposure date, a forecast object, and the
common information environment. Let $E=1$ be a specified extreme direct
experience and $E=0$ its comparison. The causal persistence profile is
\[
 \psi_h=E[B_h(1)-B_h(0)],\qquad h=0,1,\ldots.
\]
Nagel's lecture explicitly asks about stronger persistence of extreme
experiences and about direct versus indirect experience
\cite[PDF pp. 10--11]{nagel}. It does not establish universal slower
decay. Randomized exposure with consistency, independent units, complete
follow-up, and positive assignment probabilities identifies $\psi_h$ as
a difference of observed means. This total effect includes changes in
wealth and other circumstances. The results below show why even an
exactly observed, apparently well-behaved decay profile need not identify
memory persistence.

\section*{Two exponential channels and their ambiguity}
Consider a structural mean response consisting of a memory contribution
and a circumstance contribution,
\[
 \psi_h=\alpha r^h+\beta s^h,
 \qquad \alpha,\beta>0,\quad 0<r,s<1,\quad r\ne s.        \tag{1}
\]
The labels ``memory'' and ``circumstances'' refer to mechanisms, not to
which numerical rate happens to be larger. Suppose only exposure and
forecasts are observed, with no intervention that distinguishes the
channels. Residual forecast noise can have any fixed common law,
independent of randomized exposure.

\begin{theorem}
Four consecutive population responses $\psi_0,\ldots,\psi_3$ identify
the unordered rate set $\{r,s\}$ and the matching amplitudes in (1),
but do not identify which rate is the memory rate. Define
\[
 D=\psi_0\psi_2-\psi_1^2.
\]
Then $D>0$ and the two rates are the roots of
\[
 z^2-Sz+P=0,\qquad
 S=\frac{\psi_0\psi_3-\psi_1\psi_2}{D},\quad
 P=\frac{\psi_1\psi_3-\psi_2^2}{D}.                      \tag{2}
\]
After choosing which root is called $r$, its amplitude is
$\alpha=(\psi_1-s\psi_0)/(r-s)$ and $\beta=\psi_0-\alpha$.
\end{theorem}
\begin{proof}
Expanding gives $D=\alpha\beta(r-s)^2>0$. Each sequence $r^h$ and
$s^h$ satisfies the recurrence
$\psi_{h+2}=(r+s)\psi_{h+1}-rs\psi_h$; so does their linear
combination. Applying the recurrence at $h=0,1$ and solving its two
linear equations gives (2). The roots are distinct, and the equations
$\alpha+\beta=\psi_0$, $\alpha r+\beta s=\psi_1$ then give the
amplitudes. Conversely, interchange $(\alpha,r)$ and $(\beta,s)$.
Let a latent memory variable carry the first term and a latent
circumstance variable the second in one structural model, and reverse
their roles in another. Add the same residual-noise process to the
sum. These two models give exactly the same joint law of exposure
and every forecast horizon, but different memory rates. Therefore
even infinitely many forecast observations cannot distinguish the labels
in this class.
\end{proof}
The positive-amplitude and distinct-rate restrictions are essential to
this inversion. If one amplitude is zero or $r=s$, then $D=0$ and the
four equations no longer identify two rates. Small positive $D$ also
makes the inversion poorly conditioned; observing four noisy sample
means is not equivalent to knowing four population responses exactly.

\section*{A false inference from normalized total decay}
\begin{example}
Let both ordinary and extreme experiences have the same memory rate
$r=1/5$ and circumstance rate $s=9/10$. Suppose their amplitudes are
$(\alpha,\beta)=(9,1)$ and $(1,9)$, respectively. Both initial total
responses equal ten. For every $h\ge1$,
\[
 \frac{\psi_h^{\mathrm{extreme}}}{\psi_0^{\mathrm{extreme}}}
 -\frac{\psi_h^{\mathrm{ordinary}}}{\psi_0^{\mathrm{ordinary}}}
 =\frac45\{(9/10)^h-(1/5)^h\}>0.
\]
Thus extreme total responses decay strictly more slowly at every positive
horizon although the memory decay rate is identical. The difference
comes entirely from the relative channel amplitudes.
\end{example}
This counterexample is compatible with randomized exposure and exact
measurement. Holding common news fixed does not fix the circumstance
channel and hence does not prevent the example.

\section*{A channel intervention that resolves the label}
Augment the experiment with a randomized reset $Z\in\{0,1\}$ of the
post-exposure circumstance path, independent of exposure. A successful
reset places that path at the same baseline under both exposures.
Assume it neither changes memory nor affects beliefs through any other
route. These are exclusion restrictions on the reset, not consequences
of its randomization. In the two-rate model they imply
\[
 E[B_h(e,z)]=\mu_{h,z}+e\{\alpha r^h+(1-z)\beta s^h\}.    \tag{3}
\]
The term $\mu_{h,z}$ permits an exposure-independent direct reset effect.

\begin{proposition}
With independent randomized assignment of $(E,Z)$, positive probability
for all four cells, consistency and complete follow-up, the reset-cell
exposure contrast $\phi_h$ identifies $\alpha r^h$. If $\phi_0\ne0$,
then $r=\phi_1/\phi_0$ and $\alpha=\phi_0$. Subtracting $\phi_h$
from the unreset exposure contrast identifies the circumstance profile.
\end{proposition}
\begin{proof}
Randomization identifies each potential mean in (3). Subtract the
$e=0$ mean from the $e=1$ mean within $z=1$; $\mu_{h,1}$ cancels and
the circumstance term is zero. The two ratios and the unreset-minus-reset
difference then follow algebraically. No cross-world joint distribution
is needed for these mean contrasts.
\end{proof}

\section*{Ratios and the remaining empirical gap}
For any estimated profile satisfying
$|\widehat\psi_h-\psi_h|\le\varepsilon$ and
$|\widehat\psi_0-\psi_0|\le\varepsilon<|\psi_0|$,
direct subtraction of fractions gives
\[
 \left|\frac{\widehat\psi_h}{\widehat\psi_0}
            -\frac{\psi_h}{\psi_0}\right|
 \le\frac{\varepsilon}{|\psi_0|-\varepsilon}
       \left(1+\left|\frac{\psi_h}{\psi_0}\right|\right).  \tag{4}
\]
Indeed the numerator is
$(\widehat\psi_h-\psi_h)\psi_0
-\psi_h(\widehat\psi_0-\psi_0)$ and the denominator has absolute
value at least $|\psi_0|(|\psi_0|-\varepsilon)$. Near a zero initial
effect, normalized decay can therefore be unstable even with accurate
level estimates. Sign reversals should be reported in levels instead of
being hidden by a fitted positive decay rate.

The outstanding task is to find exposure and channel interventions with
credible exclusions in an actual population, handle attrition and changing
news, and test the restrictive two-rate representation. The algebra
identifies a precise ambiguity and a design that removes it under stated
conditions; it does not establish a psychological memory law from total
forecast persistence alone.

\section{Completion Estimate}
47\%

This is a post hoc, heuristic estimate for the full catalogue target.
The attempt identifies two unlabeled exponential persistence channels, demonstrates that slower total decay need not reflect slower memory decay, and supplies a reset design and ratio-error bound. Actual extreme-experience persistence still needs credible exposure and channel interventions, exclusions, attrition control and validation of the restrictive decay representation in the target population.

\begin{thebibliography}{9}
\bibitem{nagel} S. Nagel (2021), \emph{Subjective Beliefs in Asset Pricing},
Innsbruck lectures, June 2021, PDF pp. 10--11.
\url{https://bpb-us-w2.wpmucdn.com/voices.uchicago.edu/dist/f/575/files/2021/07/Innsbruck2021_slides.pdf}.
\end{thebibliography}

\end{document}
